% sections/02_related_work.tex — Section II (Related Work)
% Source: notes/paper_section_II_draft.md
% Converted Day 5 2026-05-13.

\section{Related Work}
\label{sec:related_work}

\subsection{Conditional Flow Matching and Shortcut Variants}
\label{sec:related:cfm}

Conditional Flow Matching (CFM)~\cite{lipman2023flow_matching, liu2022rectified} is a
generative modeling framework that learns a continuous-time velocity field transporting
samples from a simple source distribution (typically isotropic Gaussian) to the target
data distribution conditioned on auxiliary context. Unlike score-based diffusion models,
CFM directly parameterizes the velocity field rather than the score function, yielding
simpler training objectives and tighter connections to optimal transport. The CFM training
loss
\begin{equation}
\mathcal{L}_{\text{CFM}}(\theta) \;=\; \mathbb{E}_{t,\,x_0,\,x_1}\!\Bigl[\| v_\theta(x_t, t \mid c) - (x_1 - x_0) \|^2\Bigr],
\label{eq:cfm_loss}
\end{equation}
where $x_t = (1-t)\, x_0 + t\, x_1$ and $c$ is the conditioning signal, admits the
population minimizer $v^*(x_t, t \mid c) = \mathbb{E}_{x_1 \sim \dataD(\cdot \mid c)}[x_1
- x_0 \mid x_t, t, c]$, recovering the matched-flow velocity field. The induced sampling
distribution at $t = 1$ matches $\dataD(x_1 \mid c)$ at convergence. This recovery property
places CFM within the family of distribution-matching losses to which Theorem~\ref{thm:bound}
(Remark~\ref{rem:arch_agnostic}) applies.

A practical limitation of CFM is that high-quality sampling typically requires many Euler
integration steps. Shortcut models~\cite{frans2024shortcut} address this by additionally
conditioning the velocity network on a desired step size, enabling few-step or single-step
sampling via a self-consistency objective that ties large-step prediction to composed
small-step trajectories. The Shortcut Flow Matching variant (the head and decoder loss of
OTP-Soft, \S\ref{sec:architecture}) inherits the Lipman-style distribution-matching
structure while accelerating training-time convergence relative to vanilla CFM.

Within image generation, CFM and its variants underpin state-of-the-art models (Stable
Diffusion~3, FLUX) and achieve high sample quality on language-conditioned generation.
Applications to robot action sequences are more recent: $\pi_0$~\cite{black2024pi0} uses
vanilla flow matching for action prediction in a generalist robot policy, and several
follow-up works (including OTP-Soft in this paper) explore variants. Our finding
(\S\ref{sec:results:cfm_collapse}) is that the demonstration-supervised setting introduces
a failure mode not encountered in image-generation CFM: under language-orthogonal
demonstration distributions, the matched-flow representation cannot encode
language-discriminative action distributions because such distributions are absent in the
supervision data. This is a setting-specific failure rather than an architectural defect
of CFM as a generative framework.

\subsection{Vision-Language-Action Architectures}
\label{sec:related:vla}

Modern VLAs are built by integrating a pretrained vision-language model (VLM) with an
action head trained on robot demonstration datasets~\cite{brohan2023rt2, kim2024openvla}.
The pretrained VLM contributes broad semantic and visual knowledge from internet-scale
data; the action head learns to map VLM hidden states to robot control commands using
teleoperated demonstrations as supervision. This decomposition allows VLAs to inherit the
language-following capability of the underlying VLM while specializing to embodied control
tasks.

VLA designs differ primarily in the action head architecture and loss family.
Table~\ref{tab:vla_archs} summarizes representative architectures and their relationship
to Theorem~\ref{thm:bound}.

\begin{table}[t]
\centering
\caption{\textsc{VLA Architectures and Theorem~\ref{thm:bound} Status.}
All listed architectures use distribution-matching loss families whose population
minimizers recover $\dataD(\cdot \mid o, \ell)$; consequently all are bounded by
Theorem~\ref{thm:bound} Remark~\ref{rem:arch_agnostic}.}
\label{tab:vla_archs}
\small
\begin{tabular}{lll}
\toprule
\textbf{Architecture} & \textbf{Action head} & \textbf{Loss family} \\
\midrule
OpenVLA~\cite{kim2024openvla}      & Autoregressive tokens   & Cross-entropy (discretized) \\
OpenVLA-OFT~\cite{kim2025openvla_oft} & Parallel decoding head  & $L_1$ regression \\
RT-2~\cite{brohan2023rt2}           & Co-fine-tuned VLM tokens & Cross-entropy (discretized) \\
$\pi_0$~\cite{black2024pi0}         & Flow matching head      & Vanilla CFM \\
Octo~\cite{octo2024}                & Diffusion head          & Denoising score matching \\
DiffPolicy~\cite{chi2023diffusion}  & Conditional DDPM        & Denoising score matching \\
OTP-Soft (\S\ref{sec:architecture}) & Object-pose + decoder   & Shortcut Flow Matching \\
\bottomrule
\end{tabular}
\end{table}

The diversity of action-head choices is typically motivated by efficiency considerations
(parallel decoding versus autoregressive, continuous versus discretized actions) or by
representational flexibility (flow matching and diffusion permit multi-modal action
distributions; cross-entropy on discretized actions cannot). Comparative studies have
evaluated these design choices primarily through task success rate on benchmark suites
(LIBERO, RoboCasa, Bridge), often using paraphrase invariance as a secondary axis
interpreted as evidence of ``language understanding''~\cite{kim2025openvla_oft}.

Our work makes a different observation: under Theorem~\ref{thm:bound}
(Remark~\ref{rem:arch_agnostic}), all of the architectures in Table~\ref{tab:vla_archs}
are subject to the same bound $I(\Astar; \ell \mid o) \le I_{\dataD}(a; \ell \mid o)$,
because each uses a distribution-matching loss whose population minimizer recovers
$\dataD(\cdot \mid o, \ell)$. The architectural choice does not change this bound; it
changes the empirical character of policy behavior when the bound is tight. In
\S\ref{sec:results:oft_contrast} we discuss this in detail for the OFT contrast: OFT's
97\% SR on LIBERO-Spatial is consistent with two scenarios within the bound (low-magnitude
language preservation or entirely visual grounding), neither of which can be distinguished
by benchmark SR alone.

\subsection{Language-Grounding Probing Methodology}
\label{sec:related:probing}

Probing studies~\cite{tenney2019bert_probing, hewitt2019structural_probing} measure where
linguistic structure is encoded in the internal representations of a neural network. The
canonical setup trains a lightweight classifier on top of fixed model representations to
predict linguistic targets (part-of-speech tags, syntactic dependency structure, semantic
roles); high classifier accuracy is interpreted as evidence that the underlying
representation encodes the target. More recent variants (Tuned
Lens~\cite{belrose2023tuned_lens}, Logit Lens) probe layer-wise predictions in large
language models without auxiliary classifier training.

In language-conditioned robotics, probing has been applied to learned policies for
instruction sensitivity: e.g., measuring policy behavior under input perturbations,
attribution analyses on attention weights, or representation-level inspection of policy
hidden states. The shared assumption is that the policy is the object of inquiry ---
probing happens \emph{after} training to assess what the model learned.

Our diagnostic methodology departs from this convention in one substantive way. We retain
backbone-level probing (Gate~1, \S\ref{sec:results:diagnostic}: a Mantel test on backbone
hidden states versus language embeddings, analogous to standard NLP representation
probing), but we additionally probe the \textbf{supervision data itself} (M3,
\S\ref{sec:results:diagnostic}: a Mantel test on demonstration trajectories versus language
embeddings, prior to and independent of any specific policy). To our knowledge, this
protocol --- probing the language-trajectory correlation structure of the demonstration
dataset directly --- has not previously been applied in the language-conditioned robotics
literature. The novelty is methodological rather than technical: the Mantel test is
standard, but its application to demonstration data as a diagnostic layer separate from
model-internal probing reveals where language conditioning is preserved (backbone) versus
lost (demonstration data), prior to any policy being trained.

We also retain policy-level probing (C1, \S\ref{sec:results:diagnostic}: Mantel test on
policy latent representations versus language embeddings, analogous to standard policy
probing), so that the three-layer diagnostic (Gate~1 / M3 / C1) localizes language
information at each stage of the VLA pipeline. The relationship between these three levels
is then formalized by Theorem~\ref{thm:bound} (\S\ref{sec:theory}).
